\documentclass[12pt,a4paper]{article}
\usepackage[utf8]{inputenc}
\usepackage[T1]{fontenc}
\usepackage{lmodern}
\usepackage{amsmath,amssymb}
\usepackage{siunitx}
\usepackage{geometry}
\usepackage{hyperref}
\usepackage{enumitem}
\usepackage{tikz}
\geometry{margin=2.2cm}

\title{Lecture 7 -- Dual-Slope ADC}
\author{Readable Assignment Set}
\date{}

\begin{document}
\maketitle
\tableofcontents

\section{Assignments}
\subsection{G1}
Dual-slope ADC analog section with integration interval $T_1=\SI{20}{ms}$,
$0\le e_{in}\le\SI{10}{V}$,
$E_R=\SI{-2.5}{V}$,
$R=\SI{100}{k\Omega}$,
$C=\SI{150}{nF}$,
and comparator/op-amp saturation data $V_{SAT}^+=-V_{SAT}^-=\SI{15}{V}$.

For ideal A1 and A2 in Questions 1 and 2:
\begin{enumerate}
\item With $e_{in}=\SI{10}{V}$, find integrator output $e_1$ after $T_1$.
\item If measured $T_2=\SI{56}{ms}$, find mean input $\bar e_{in}$ during $T_1$. Sketch conversion process for $e_1$ and $e_2$ with slopes and levels.
\end{enumerate}
For Question 3 include A1 non-idealities: $V_{OS}=\pm\SI{5}{mV}$, $I_B=\SI{0.5}{\micro A}$ (npn), $I_{OS}=0$.
\begin{enumerate}[resume]
\item Find maximum percentage conversion error for $e_{in}=\bar e_{in}$ from Q2 (or \SI{7}{V} if Q2 unavailable).
\end{enumerate}

\subsection{G2}
Design an 18-bit dual-slope ADC using Figure 12.50, with $V_{ref}=\SI{5}{V}$ and \SI{5}{MHz} clock. Input $v_s$ ranges \SIrange{0}{10}{V}. Fixed interval $T_1$ corresponds to count $2^n$.
Find maximum conversion time. If max $V_{peak}=\SI{10}{V}$ and $R=\SI{1}{k\Omega}$, determine needed $C$. If $R$ increases by a few percent with aging, discuss converter accuracy and whether max $V_{peak}$ increases or decreases.

\section{Hints}
\begin{itemize}
\item Dual-slope principle: integrate unknown for fixed time $T_1$, then deintegrate with known reference until output returns to zero.
\item Integrator equation: $\Delta v = -\dfrac{1}{RC}\int v_{in}(t)\,dt$.
\item For constant signals, $T_2=\dfrac{|\bar e_{in}|}{|E_R|}T_1$.
\item Clock count is proportional to measured deintegration time, giving excellent noise averaging and component-insensitive gain ratio.
\end{itemize}

\section{Step-by-step solutions}
\subsection{G1.1}
\[
\frac{de_1}{dt}=-\frac{e_{in}}{RC}=-\frac{10}{(100k)(150n)}=-\SI{666.7}{V/s}.
\]
After $T_1=\SI{20}{ms}$:
\[
\Delta e_1 = \frac{de_1}{dt}T_1=-666.7\cdot0.02=-\SI{13.33}{V}.
\]
So $e_1(T_1)=\SI{-13.33}{V}$ (assuming start at 0 V and no saturation limit reached first).

\subsection{G1.2}
For deintegration with $E_R=-\SI{2.5}{V}$ and opposite slope sign:
\[
T_2=\frac{|\bar e_{in}|}{|E_R|}T_1
\Rightarrow
\bar e_{in}=|E_R|\frac{T_2}{T_1}=2.5\cdot\frac{56}{20}=\SI{7.0}{V}.
\]
Hence mean input is \SI{7.0}{V}. During $T_1$, $e_1$ ramps linearly downward; during $T_2$, it ramps upward back to 0. Comparator output toggles when crossing 0.

\subsection{G1.3}
Include offset and bias currents as equivalent input-error voltage to integrator:
\[
e_{err}=V_{OS}+I_B R_{eq}.
\]
Relative conversion error (first order):
\[
\varepsilon_\%\approx\frac{|e_{err}|}{|e_{in}|}\cdot100\%.
\]
Using $e_{in}=\SI{7}{V}$ and inserting worst-case $V_{OS}=\SI{5}{mV}$ plus bias term gives a small but finite ppm/percent-level error; evaluate with exact $R_{eq}$ from the input network.

\subsection{G2}
\begin{enumerate}
\item For 18-bit, $n=18$ and fixed run-up count is $2^{18}=262144$ clocks.
\item With $f_{clk}=\SI{5}{MHz}$:
\[
T_1=\frac{2^{18}}{5\times10^6}=\SI{52.43}{ms}.
\]
\item Maximum deintegration occurs at full-scale $v_s=\SI{10}{V}$ with $V_{ref}=\SI{5}{V}$:
\[
T_{2,max}=\frac{10}{5}T_1=2T_1=\SI{104.86}{ms}.
\]
\item Maximum conversion time:
\[
T_{conv,max}=T_1+T_{2,max}=\SI{157.29}{ms}.
\]
\item Integrator ramp limit relation:
\[
\left|\frac{dv}{dt}\right|=\frac{V_{peak}}{RC} \Rightarrow C=\frac{V_{peak}}{R\cdot |dv/dt|_{design}}.
\]
Use the allowed internal ramp/swing target from design constraints to set $|dv/dt|_{design}$ and solve $C$.
\item If $R$ drifts upward, both integration and deintegration slopes scale similarly, so ratio-based conversion accuracy ideally remains nearly unchanged. However, dynamic range/headroom and maximum ramp amplitude trend downward (slower slope), so effective peak excursion decreases for fixed timing.
\end{enumerate}

\subsection{Timing sketch}
\begin{center}
\begin{tikzpicture}[xscale=0.95,yscale=0.7]
\draw[->] (0,0) -- (12,0) node[right] {$t$};
\draw[->] (0,-3) -- (0,3);
\draw[thick] (0,0) -- (4,-2.2) -- (9,0);
\node at (2.1,-2.7) {$T_1$};
\node at (6.5,-2.7) {$T_2$};
\node at (9.6,0.3) {$e_1$};
\draw[thick] (0,2.2) -- (9,2.2) -- (9,-2.2) -- (11,-2.2);
\node at (11.2,-2.2) {$e_2$};
\end{tikzpicture}
\end{center}

\end{document}
